\documentclass[10pt]{article}
\usepackage[margin=0.7in]{geometry}
\usepackage{booktabs,graphicx,amsmath,hyperref}
\begin{document}
\section*{Campaign findings: extensive\_20260907}
Completed 770/770 configured training runs and 11/11 reference kernels.

Configuration extensive\_20260907: 11 PDEs, widths [4, 16, 64, 256, 1024, 4096, 16384], 10 seeds, 3000 full-batch GD steps. All parameters of the biased two-layer tanh network are trained in float64; loss is ||P-y||\textasciicircum{}2 and eta=0.2/lambda\_max(K0).

\paragraph{Widest completed width per PDE.}
Seed medians; loss columns are $L_T/L_0$. Drift is the maximum recorded relative operator-norm change.
\begin{center}\small\begin{tabular}{lrrrrrr}\toprule
PDE & Width & Seeds & Relative gap & GD loss & Frozen loss & Drift \\
\midrule
Poisson (1D) & 16384 & 10/10 & 0.0006 & 4.62e-06 & 5.63e-06 & 0.0155 \\
Variable elliptic (2D) & 16384 & 10/10 & 0.00327 & 6.91e-08 & 6.87e-08 & 0.0048 \\
Heat (1+1D) & 16384 & 10/10 & 2.32e-06 & 0.000189 & 0.00019 & 0.000839 \\
Transport (1+1D) & 16384 & 10/10 & 6.49e-07 & 6.22e-05 & 6.33e-05 & 0.000367 \\
Wave (1+1D) & 16384 & 10/10 & 1.88e-07 & 0.000207 & 0.000206 & 0.00175 \\
Biharmonic (2D) & 16384 & 10/10 & 1.96e-07 & 0.00522 & 0.00523 & 0.00486 \\
Stokes (2D, vector) & 16384 & 10/10 & 0.00017 & 0.000448 & 0.00045 & 0.000595 \\
Weak Poisson (1D) & 16384 & 10/10 & 0.000131 & 9.77e-05 & 9.75e-05 & 0.000722 \\
Nonlocal diffusion (1D) & 16384 & 10/10 & 2.86e-05 & 2.29e-05 & 2.29e-05 & 0.000432 \\
Integro-differential (1D) & 16384 & 10/10 & 3.16e-05 & 0.000124 & 0.000124 & 0.000474 \\
Caputo diffusion (1+1D) & 16384 & 10/10 & 2.13e-06 & 0.00128 & 0.00129 & 0.00213 \\
\bottomrule\end{tabular}\end{center}
Of the 770 completed runs, 770 reduce training loss and 767 reduce it by at least a factor of ten within the configured budget. These counts include all PDEs and widths and are descriptive, not a test of a common convergence rate.

\begin{figure}[ht]\centering\includegraphics[width=\linewidth]{\detokenize{figures/conditioning_loss_relationship.pdf}}
\caption{Conditioning and final loss, separated by PDE. Circles show nonlinear GD, plus signs show the initial frozen kernel, marker size denotes width, and color denotes checkpoint kernel drift. Each seed is shown; no pooled fit is made.}\end{figure}
\paragraph{Interpretation limits.}\begin{itemize}
\item These summaries describe completed numerical runs. Unfinished cells remain in the coverage denominator; missing work is not a failed convergence experiment.
\item All registered PDE designs have analytic independence routes under the paper's assumptions. QMC estimates the gap; an unresolved or nonpositive numerical eigenvalue does not refute limiting positive definiteness.
\item The spectral learning-rate rule and tested widths do not verify the theorem's deterministic derivative bounds, sufficient width inequality, or all-iteration guarantee.
\item Positive definiteness alone does not predict finite-budget loss reduction: the relative gap, initial modal energy, finite width, and kernel drift also matter. The plot is descriptive and does not fit a pooled causal relationship across PDEs.
\item The largest width reported for a PDE is the largest width with accepted completed data; its seed count can be incomplete. Full per-PDE/per-width coverage is in findings\_by\_problem\_width.csv.
\item Kernel drift and gap preservation are checked at recorded checkpoints. Kernel drift is normalized by the largest eigenvalue and can be small while exceeding the smallest eigenvalue.
\item Weak, nonlocal, and Caputo experiments optimize their stored finite quadrature constraints. Training loss does not establish continuum solution accuracy or generalization.
\item The existing main-paper conditioning run uses a different output-bias and loss convention. Its numerical values are not combined with this campaign.
\end{itemize}
\noindent Configuration SHA-256: \texttt{93f8907b121f8baa0c2a20a461b30df6e42e34ded852de0ca308593bc753a689}.
\end{document}
